\documentclass[12pt, titlepage]{article}

\usepackage{graphicx} % Allows for images
\usepackage{wrapfig} % Allows for text wrapping around images
\usepackage{hyperref} % Allows for hyperlinks
\usepackage{siunitx} % Allows for SI units
\usepackage{amsmath} % Allows for math
\usepackage{enumerate} % Allows for custom enumerations
\usepackage{xcolor} % Allows for custom colors
\usepackage{tikz, tcolorbox} % Allows for custom boxes
\usepackage{microtype} % Allows for better text alignment
\usepackage{listings} % Allows for code listings
\usepackage{booktabs} % Allows for better tables
\usepackage{float} % Allows for better figure placement
\usepackage{geometry} % Allows for custom page geometry
\usepackage{fancyhdr} % Allows for custom headers and footers
\usepackage{caption} % Allows for custom captions

\fancypagestyle{myarticlestyle}{
    \fancyhf{} % Clear header and footer
    \fancyhead[L]{\leftmark} % Section name on the left
    \fancyhead[R]{\thepage} % Page number on the right
}
\fancypagestyle{tocstyle}{
    \fancyhf{} % Clear header and footer
    \fancyhead[L]{CONTENTS} % Section name on the left
    \fancyhead[R]{\thepage} % Page number on the right
}

\pagestyle{myarticlestyle}
\setlength{\headheight}{15pt}

\begin{document}
\tableofcontents
\listoffigures
\listoftables
\thispagestyle{tocstyle}
\newpage
\section{Objectives}
This laboratory experiment aims to determine the principal stresses at multiple
points within a cantilevered tube subjected to distinct loading conditions:

\subsection{Bending and Torsional Load Analysis}
\begin{itemize}
    \item Investigate and analyze the principal stresses occurring in the
      cantilevered tube when subjected to combined bending and torsional loads.
    \item Determine the stress distribution and variations at various points
      along the length of the tube under the influence of these specific loads.
\end{itemize}

\subsection{Bending, Torsional Load, and Internal Pressure Analysis}
\begin{itemize}
    \item Examine the effect of an internal pressure in conjunction with
      bending and torsional loads on the cantilevered tube. 
    \item Identify and
      evaluate the changes in principal stresses at different critical
      locations within the tube due to the combined influence of internal
      pressure, bending, and torsional loading.
\end{itemize}

These objectives aim to provide a comprehensive understanding of the stress
behavior and distribution within the cantilevered tube under various loading
conditions, specifically focusing on the determination of principal stresses.
The results obtained will aid in characterizing the mechanical response of the
structure and provide insights for practical engineering applications and
design considerations.
\newpage
\section{Procedure}
This section outlines the procedure for the laboratory experiment, including
the equipment and materials used, the testing procedure, and the data
collection process. The sketch of the tabular structure with dimensions and
strain gauge locations is shown in Figure INS in the Appendix.
\subsection{Preliminary Set-Up Procedure}
\begin{enumerate}
    \item Check the specimen dimensions and strain gauge location.
    \item Check lead wires and identify each gauge with the number of the
      circuit in the Datascan Analog Measurement Processor.
    \item Hang the loading pan on the side arm of the tube at a distance of
      457.2 mm (18 in.) from the tube axis.
\end{enumerate}

\subsection{Testing Procedure}
\begin{enumerate}
    \item Check and record the initial values for each of the six strain gauges
      in the unloaded configuration.
    \item Loading A:
    \begin{enumerate}
        \item Apply loads in six equal increments of 20 lbs (total 120 lbs) for
          the aluminum specimen and increments of 40 lbs (total 240 lbs) for
          the steel specimen.
        \item Record the readings of strain for each gauge and at each load.
          Unload the specimen.
    \end{enumerate}
    \item Loading B:
    \begin{enumerate}
        \item Supply and maintain a constant air pressure in the specimens (100
          psi in the aluminum specimen, and 200 psi in the steel specimen), as
          indicated by the pressure gauge.
        \item Apply loads in six equal increments of 20 lbs (total 120 lbs) for
          the aluminum specimen and increments of 40 lbs (total 240 lbs) for
          the steel specimen.
        \item Record the readings of strain for each gauge and at each load.
          Unload the specimen.
    \end{enumerate}
\end{enumerate}
\newpage
\section{Results}

\begin{table}[H]
\begin{tabular}{|c|c|c|c|c|c|c|} \hline
\multicolumn{7}{|c|}{Strain gauge readings (10$^{-6}$)} \\ \hline
Load (lb) & 1 & 2 & 3 & 4 & 5 & 6 \\ \hline
0 & 0.0 & 0.0 & 0.0 & 0.0 & 0.0 & 0.0 \\ \hline
20 & & & & & & \\ \hline
40 & & & & & & \\ \hline
60 & & & & & & \\ \hline
80 & & & & & & \\ \hline
100 & & & & & & \\ \hline
120 & & & & & & \\ \hline
\end{tabular}
\caption{Strain gauge readings for aluminum specimen}
\end{table}
\newpage
\section{Discussion}
\newpage
\section{Conclusions}
\newpage
\appendix
\section{Appendix}
\end{document}
